\documentclass[11pt]{article}
\usepackage{amsmath,amssymb,amsthm}
\usepackage[margin=1in]{geometry}
\usepackage{booktabs}
\usepackage[expansion=false]{microtype}
\usepackage[colorlinks=true,linkcolor=blue,urlcolor=blue]{hyperref}

\newtheorem{proposition}{Proposition}
\newtheorem{remark}{Remark}
\newtheorem{definition}{Definition}

\newcommand{\R}{\mathbb{R}}
\newcommand{\E}{\mathbb{E}}
\newcommand{\sg}{\operatorname{sg}}
\newcommand{\softmax}{\operatorname{softmax}}
\newcommand{\1}{\mathbf{1}}
\DeclareMathOperator{\Cat}{Cat}

\title{\bf The Cue-Triggered Lick-Timing Agent:\\ Network, Task and Learning Rule}
\author{neuralgeom / \texttt{tasks.timing}}
\date{\today}

\begin{document}
\maketitle

\noindent
This note writes out, without abbreviation, what one training update computes. Every
symbol is defined at first use. The reference implementation is
\texttt{generator.py}, \texttt{rl.py} and \texttt{variants.py} under
\texttt{neuralgeom/tasks/timing/}.

\tableofcontents

\section{Index conventions and dimensions}\label{sec:notation}

\begin{definition}[Indices]
Three indices run throughout.
\begin{itemize}
  \item $b \in \{1,\dots,B\}$ indexes the \emph{environment copy}. $B = 16$ independent
        trial generators run in lockstep. They share one set of network weights but keep
        separate trial histories, outcome histories and delay schedules; they exist only to
        average the gradient, and are never coupled. All equations below are written for a
        single $b$ and the index is suppressed except where the batch is summed over.
  \item $t \in \{0,1,\dots,T-1\}$ indexes the \emph{time step within one episode}. One
        episode is one trial. The physical step size is $\mathrm{d}t = 20$\,ms, so
        $t$ counts $20$\,ms bins. $T$ is the length of the longest trial in the batch;
        shorter trials are padded (\S\ref{sec:mask}).
  \item $i \in \{1,\dots,I\}$ indexes the \emph{gradient update}. One update consumes one
        batch of $B$ episodes. $I$ is the \texttt{--steps} argument.
\end{itemize}
\end{definition}

\begin{definition}[Dimensions]
$N = 128$ is the number of recurrent units; $d_u = 5$ is the observation dimension
(\S\ref{sec:obs}); $|\mathcal{A}| = 2$ is the number of actions.
\end{definition}

\begin{definition}[Stop-gradient]
$\sg[\,\cdot\,]$ denotes the identity in the forward pass and the zero map in the backward
pass: $\sg[x] = x$ and $\partial\,\sg[x]/\partial\theta = 0$. It is PyTorch's
\texttt{.detach()}. Where it appears, the enclosed quantity is treated as a constant by the
optimiser.
\end{definition}

\section{The observation}\label{sec:obs}

\begin{definition}[Observation vector]
At step $t$ the agent receives $u_t \in \R^{5}$,
\begin{equation}
  u_t \;=\; \bigl(\, c_t,\; \rho_{-1},\; \alpha_{-1},\; \varsigma_{-1},\; \phi_{-1} \,\bigr)^{\!\top},
\end{equation}
with components
\begin{align}
  c_t        &= \1\bigl[\,0 \le t - t_{\mathrm{cue}} < n_{\mathrm{cue}}\,\bigr], \\
  \rho_{-1}  &= \textstyle\sum_{t} r_t^{\,\mathrm{prev}}, \\
  \alpha_{-1}    &= \1[\text{the previous trial contained a decisive lick}], \\
  \varsigma_{-1} &= \1[\text{the previous trial was rewarded}], \\
  \phi_{-1}      &= \text{first-lick latency on the previous trial, in seconds ($0$ if none)}.
\end{align}
In words: $c_t$ is the \emph{cue channel}, equal to $1$ exactly while the cue is audible;
$\rho_{-1}$ is the total reward accumulated over the whole of the previous trial; and
$\alpha_{-1}, \varsigma_{-1}, \phi_{-1}$ report what the agent did on that trial and how it
ended.

Here $t_{\mathrm{cue}}$ is the step of cue onset ($t_{\mathrm{cue}} = \varnothing$ before it,
in which case $c_t = 0$), and $n_{\mathrm{cue}} = \lceil 0.6\,\mathrm{s}/\mathrm{d}t \rceil = 30$
steps.
\end{definition}

Two properties of $u_t$ matter for what follows.

\begin{remark}[The cue is a signal, not a phase]
$c_t$ depends only on time since cue onset. It is not tied to the reward contingency: if the
required delay is shorter than the cue duration, the cue is still audible during the rewarded
window. The agent is never told which phase the trial is in.
\end{remark}

\begin{remark}[The history channels are per-trial constants]
$\rho_{-1},\alpha_{-1},\varsigma_{-1},\phi_{-1}$ are written once, when the previous trial
ends, and held fixed for every step of the current trial. They are a static offset, not an
event. This is what allows an episode to be a single trial: cross-trial information reaches
the network through the input rather than through a persistent hidden state.
\end{remark}

\section{The network}\label{sec:net}

The recurrent core is a leaky (continuous-time) $\tanh$ network. Let $h_t \in \R^{N}$ be the
hidden state after step $t$.

\begin{definition}[State equation]
\begin{align}
  a_t^{\mathrm{pre}} &= W_{\mathrm{rec}}\, h_{t-1} \;+\; W_{\mathrm{in}}\, u_t \;+\; b_{\mathrm{in}}
    \;+\; \sqrt{\tfrac{2}{\alpha}}\,\sigma_{\!\xi}\,\xi_t,
    \qquad \xi_t \sim \mathcal{N}(0, I_N), \label{eq:pre}\\[2pt]
  h_t &= (1-\alpha)\, h_{t-1} \;+\; \alpha\, \tanh\!\bigl(a_t^{\mathrm{pre}}\bigr),
    \qquad h_{-1} = h_0 = 0. \label{eq:state}
\end{align}
\end{definition}

\noindent Symbols in \eqref{eq:pre}--\eqref{eq:state}:
\begin{itemize}
  \item $h_t \in \R^{N}$ --- the \emph{hidden state} (lowercase $h$ throughout). Not to be
        confused with $\mathcal{H}$, the entropy (\S\ref{sec:policy}).
  \item $a_t^{\mathrm{pre}} \in \R^{N}$ --- the pre-activation (total synaptic input).
  \item $W_{\mathrm{rec}} \in \R^{N\times N}$ --- recurrent weights, no bias. Initialised
        $(W_{\mathrm{rec}})_{jk} \sim \mathcal{N}(0, g^2/N)$ with $g=1$.
  \item $W_{\mathrm{in}} \in \R^{N\times d_u}$, $b_{\mathrm{in}} \in \R^{N}$ --- input weights
        and bias; $(W_{\mathrm{in}})_{jk}\sim\mathcal{N}(0,1/d_u)$, $b_{\mathrm{in}} = 0$ at
        initialisation.
  \item $\alpha = \mathrm{d}t/\tau = 20/100 = 0.2$ --- the \emph{leak}, the discretisation of
        $\tau \dot{h} = -h + \tanh(\cdot)$ by forward Euler. $1/\alpha = 5$ steps is the
        intrinsic memory time constant.
  \item $\sigma_{\!\xi} = 0.05$ --- private noise SD. The factor $\sqrt{2/\alpha}$ makes the
        stationary variance of the noise-driven state independent of $\mathrm{d}t$, so
        changing $\mathrm{d}t$ does not silently change the effective noise level.
  \item $\xi_t$ --- i.i.d.\ standard normal, resampled every step and every environment copy.
        Present in training mode only; \texttt{model.eval()} sets $\sigma_{\!\xi}=0$ for
        deterministic analysis.
\end{itemize}

\begin{definition}[Heads]
Two affine maps read out $h_t$:
\begin{align}
  z_t &= W_{\pi}\, h_t + b_{\pi} \;\in\; \R^{2}
    && \text{policy logits,} \quad W_{\pi}\in\R^{2\times N},\ b_\pi\in\R^2, \label{eq:logits}\\
  V_t &= w_v^{\top} h_t + b_v \;\in\; \R
    && \text{state value,} \quad w_v\in\R^{N},\ b_v\in\R. \label{eq:value}
\end{align}
At initialisation $W_\pi$ is scaled by $0.1$ and $b_\pi = 0$, so the policy starts close to
indifferent.
\end{definition}

The full parameter vector optimised is
\begin{equation}
  \theta \;=\; \bigl\{\, W_{\mathrm{in}},\, b_{\mathrm{in}},\, W_{\mathrm{rec}},\,
                        W_{\pi},\, b_{\pi},\, w_v,\, b_v \,\bigr\}.
\end{equation}
$h_0$ is fixed at $0$ and is not trained.

\section{The policy, the logit gap, and the entropy}\label{sec:policy}

\begin{definition}[Action distribution]
Write $z_t = (z_t^{\mathrm{wait}}, z_t^{\mathrm{lick}})^{\!\top}$. The policy is the softmax
over these two logits,
\begin{equation}
  \pi_t \;=\; \softmax(z_t) \;\in\; \Delta^{1},
  \qquad
  \pi_t(a) = \frac{e^{z_t^{a}}}{e^{z_t^{\mathrm{wait}}} + e^{z_t^{\mathrm{lick}}}},
  \quad a \in \{\mathrm{wait},\mathrm{lick}\}.
\end{equation}
The emitted action is sampled,
\begin{equation}
  A_t \sim \Cat(\pi_t), \qquad A_t \in \{0,1\}, \quad A_t = 1 \equiv \mathrm{lick}.
\end{equation}
\end{definition}

\begin{definition}[Logit gap $\Delta_t$ and lick probability $p_t$]
Because there are exactly two actions, the entire policy is one scalar. Define the
\emph{logit gap}
\begin{equation}
  \boxed{\;\Delta_t \;=\; z_t^{\mathrm{lick}} - z_t^{\mathrm{wait}}
    \;=\; \underbrace{\bigl(W_{\pi,\mathrm{lick}} - W_{\pi,\mathrm{wait}}\bigr)}_{=:\;w_\Delta \,\in\, \R^{N}}
      {}^{\!\top} h_t \;+\; \bigl(b_{\pi}^{\mathrm{lick}} - b_{\pi}^{\mathrm{wait}}\bigr)\;}
  \label{eq:delta}
\end{equation}
where $W_{\pi,\mathrm{lick}}, W_{\pi,\mathrm{wait}} \in \R^{N}$ are the two rows of $W_\pi$
and $w_\Delta \in \R^{N}$ is their difference --- a single direction in state space, the
\emph{policy readout direction}. Then
\begin{equation}
  p_t \;:=\; \pi_t(\mathrm{lick}) \;=\; \sigma(\Delta_t), \qquad
  \sigma(x) = \frac{1}{1+e^{-x}}.
\end{equation}
$p_t$ is the per-step lick probability. Because it is evaluated afresh at every $20$\,ms bin
and the first lick ends the decision, $p_t$ is a discrete-time \emph{hazard rate} and the
lick time is its first passage.
\end{definition}

\begin{definition}[Entropy $\mathcal{H}_t$]
The Shannon entropy of the action distribution at step $t$, in nats:
\begin{equation}
  \mathcal{H}_t \;=\; -\!\!\sum_{a \in \mathcal{A}}\! \pi_t(a)\,\log \pi_t(a)
   \;=\; -\,p_t \log p_t - (1-p_t)\log(1-p_t) \;=:\; \mathcal{H}(\Delta_t).
  \label{eq:entropy}
\end{equation}
$\mathcal{H}_t \in [0, \log 2]$, with $\log 2 = 0.693$ at $\Delta_t = 0$ (a fair coin) and
$\mathcal{H}_t \to 0$ as $|\Delta_t| \to \infty$ (a deterministic policy). \emph{$\mathcal{H}$
is a property of the output distribution, not of the hidden state $h_t$}; by \eqref{eq:delta}
it depends on $h_t$ only through the scalar projection $w_\Delta^\top h_t$.
\end{definition}

\begin{definition}[Log-likelihood of the taken action]
\begin{equation}
  \ell_t \;=\; \log \pi_t(A_t)
   \;=\; A_t \log p_t + (1-A_t)\log(1-p_t).
\end{equation}
This is the only place the sampled action enters the gradient.
\end{definition}

\section{The environment}\label{sec:env}

The generator is a five-phase machine. Let $\mathcal{P}_t \in \{\textsc{iti}, \textsc{waiting},
\textsc{eligible}, \textsc{post}, \textsc{no\_cue}\}$ be the phase \emph{before} step $t$ is
taken.

\begin{definition}[Accepted lick]
A lick action is accepted only outside a refractory period,
\begin{equation}
  L_t \;=\; A_t \cdot \1\!\left[\, t - t_{\mathrm{last}} \ge n_{\mathrm{ref}} \,\right],
  \qquad n_{\mathrm{ref}} = \lceil 0.05\,\mathrm{s}/\mathrm{d}t\rceil = 3 \text{ steps},
\end{equation}
$t_{\mathrm{last}}$ being the step of the last accepted lick. A held-down action is one lick,
not one per bin. Only $L_t$, never $A_t$, has consequences in the environment.
\end{definition}

\begin{definition}[Reward]
The scalar reward paid at step $t$ is the sum of a step cost and the phase-dependent events it
triggers:
\begin{equation}
  r_t \;=\; \underbrace{k_{\mathrm{time}}\,(1 + \vartheta_t)}_{\text{step cost}}
  \;+\;
  \begin{cases}
    k_{\mathrm{iti}}\, L_t + k_{\mathrm{to}}\,\1[\,n^{\mathrm{iti}}_t \ge n_{\mathrm{to}}\,]
      & \mathcal{P}_t = \textsc{iti},\\[4pt]
    k_{\mathrm{nc}}\, L_t
      & \mathcal{P}_t = \textsc{no\_cue},\\[4pt]
    \gamma_t k_{\mathrm{early}}\, L_t
      & \mathcal{P}_t = \textsc{waiting},\\[4pt]
    \gamma_t R_{\max}\, L_t \;+\; k_{\mathrm{miss}}\,\1[\text{window expired}]
      & \mathcal{P}_t = \textsc{eligible},\\[4pt]
    0 & \mathcal{P}_t = \textsc{post}.
  \end{cases}
  \label{eq:reward}
\end{equation}
\end{definition}

\noindent with
\begin{itemize}
  \item $\varepsilon_t = (t - t_{\mathrm{cue}})\,\mathrm{d}t$ --- seconds elapsed since cue
        onset. The delay is \emph{always} measured from cue onset.
  \item $\gamma_t > 0$ --- the \emph{value gain}, an across-trial quantity:
        $\gamma_t = \mathrm{clip}\bigl(e^{\,g\,(\rho_{\mathrm{ref}} - \rho_t)}\bigr)$ where
        $\rho_t$ is the rewarded fraction over the last 100 trials. It multiplies both the
        water and the early-lick penalty, so a run of failures makes the opportunity a trial
        represents worth more, and a run of successes worth less.
  \item $n^{\mathrm{iti}}_t$ --- steps elapsed in the current stop-licking period;
        $n_{\mathrm{to}} = 20\,\mathrm{s}/\mathrm{d}t$ is the timeout.
  \item constants $k_\bullet$, $R_{\max}$, $\kappa$ as tabulated in \S\ref{sec:consts}.
\end{itemize}

\begin{definition}[The restart rule]
In phase \textsc{iti} the remaining stop-licking counter $n^{\mathrm{rem}}$ evolves as
\begin{equation}
  n^{\mathrm{rem}} \;\leftarrow\;
  \begin{cases}
    \text{fresh draw from } \mathrm{Exp}(\mu)\big|_{[\mathrm{lo},\mathrm{hi}]}
      & \text{if } L_t = 1 \text{ and } \texttt{iti\_restart\_on\_lick},\\
    n^{\mathrm{rem}} - 1 & \text{otherwise,}
  \end{cases}
  \label{eq:restart}
\end{equation}
and the cue fires when $n^{\mathrm{rem}}$ reaches $0$. The Boolean flag in \eqref{eq:restart}
is the difference between variants F/G/I and H/J.
\end{definition}

\section{Padding and the live mask}\label{sec:mask}

Trials in a batch end at different steps. Let $D_t \in \{0,1\}$ be $1$ once the episode has
terminated. Define the \emph{live mask}
\begin{equation}
  m_t \;=\; 1 - D_{t-1}, \qquad m_0 = 1,
\end{equation}
so $m_t = 1$ exactly when the trial was still running as step $t$ was taken. After termination
the observation is frozen and the network is still stepped --- $h_t$ keeps evolving --- but
$r_t = 0$ and $m_t = 0$, so those steps contribute nothing to any sum below.

\begin{remark}
Every average in \S\ref{sec:update} is taken over live steps only, with normaliser
\begin{equation}
  \mathcal{N} \;=\; \sum_{t=0}^{T-1}\sum_{b=1}^{B} m_{t,b}.
\end{equation}
Omitting the mask would let padding drag every statistic toward zero, in proportion to how
much the trial lengths differ.
\end{remark}

\section{One update}\label{sec:update}

Fix the update index $i$.

\paragraph{Step 1: roll out.} For $t = 0,\dots,T-1$ and each $b$, run \eqref{eq:state},
\eqref{eq:logits}, \eqref{eq:value}, sample $A_t \sim \Cat(\pi_t)$, and step the environment,
storing $\ell_t$, $\mathcal{H}_t$, $V_t$, $r_t$, $m_t$, and the ITI indicator
\begin{equation}
  c^{\textsc{iti}}_t \;=\; \1[\,\mathcal{P}_t = \textsc{iti}\,] \;\in\;\{0,1\}.
\end{equation}

\paragraph{Step 2: shape the reward.}
\begin{equation}
  \tilde r_t \;=\; c_{\mathrm{scale}}\, r_t \;-\; \lambda\, \sg\!\left[\,p_t\,\right] \, c^{\textsc{iti}}_t .
  \label{eq:shaped}
\end{equation}
$c_{\mathrm{scale}} = 0.1$ rescales the environment's rewards once, globally, so that the
learning rate need not be retuned whenever a penalty changes. $\lambda \ge 0$ is
\texttt{iti\_readout\_penalty}: it charges the \emph{probability} of licking at every ITI step,
not only at steps where a lick occurred. The stop-gradient is essential --- without it the
optimiser could reduce the loss by driving $p_t$ down directly, bypassing the return, and
\eqref{eq:shaped} would no longer be a reward.

\paragraph{Step 3: returns.} With discount $\gamma$, computed backwards from $R_T = 0$:
\begin{equation}
  R_t \;=\; \tilde r_t \;+\; \gamma\, m_t\, R_{t+1}.
  \label{eq:return}
\end{equation}
$\gamma = 1$ in all runs here, so $R_t$ is the true undiscounted return to the end of the
trial. (Temporal discounting of the \emph{water} is applied inside the environment through
$e^{-\kappa\varepsilon_t}$ in \eqref{eq:reward}, which is a statement about the task's
economics, not about credit assignment.)

\paragraph{Step 4: advantage.} The value head serves as a baseline:
\begin{equation}
  \hat A_t \;=\; \sg\!\left[\, R_t - V_t \,\right],
\end{equation}
then standardised over live steps,
\begin{equation}
  \bar A_t \;=\; \frac{\hat A_t - \mu}{s},
  \qquad
  \mu = \frac{1}{\mathcal{N}}\sum_{t,b} m\,\hat A,
  \qquad
  s = \sqrt{\frac{1}{\mathcal{N}}\sum_{t,b} m\,(\hat A - \mu)^2} \;\vee\; 10^{-6}.
\end{equation}

\paragraph{Step 5: the loss.}
\begin{align}
  \mathcal{L}_{\pi} &= -\frac{1}{\mathcal{N}} \sum_{t,b} m_t\, \ell_t\, \bar A_t
    && \text{REINFORCE (policy-gradient) term,} \label{eq:lpi}\\
  \mathcal{L}_{V} &= \frac{1}{\mathcal{N}} \sum_{t,b} m_t\, \bigl(V_t - R_t\bigr)^{2}
    && \text{value regression,}\\
  \bar{\mathcal{H}} &= \frac{1}{\mathcal{N}} \sum_{t,b} m_t\, \mathcal{H}_t
    && \text{mean policy entropy,}\\[4pt]
  \mathcal{L} &= \mathcal{L}_{\pi} \;+\; c_V\,\mathcal{L}_{V} \;-\; \beta_i\, \bar{\mathcal{H}},
  && c_V = 0.5. \label{eq:loss}
\end{align}
The coefficient $\beta_i$ is annealed linearly in the update index,
\begin{equation}
  \beta_i \;=\; \beta_0 + (\beta_{\!f} - \beta_0)\,\frac{i}{I},
  \qquad \beta_0 = 0.05,\ \ \beta_{\!f} = 0.03 .
  \label{eq:beta}
\end{equation}

\paragraph{Step 6: the parameter step.}
\begin{equation}
  g \;=\; \nabla_\theta \mathcal{L},
  \qquad
  g \;\leftarrow\; g \cdot \min\!\left(1, \frac{G}{\lVert g \rVert_2}\right),\ \ G = 1,
  \qquad
  \theta \;\leftarrow\; \mathrm{Adam}\bigl(\theta,\, g,\, \eta\bigr),\ \ \eta = 4\times10^{-3}.
\end{equation}

\begin{remark}[What carries credit backwards]
$\nabla_\theta \mathcal{L}$ is taken by backpropagation through time through the full
recurrence \eqref{eq:state}: the gradient of $\ell_t$ reaches $W_{\mathrm{rec}}$ at every
earlier step. There is no eligibility trace and no truncation. Since $\hat A_t$ is
stop-gradiented, the value head is trained \emph{only} by $\mathcal{L}_V$; the policy sees
$V_t$ as a constant baseline.
\end{remark}

\section{What the entropy term does to the network}\label{sec:ent}

\begin{proposition}[The entropy bonus acts on one direction, and its force vanishes]
\label{prop:ent}
Let $\mathcal{H}(\Delta)$ be as in \eqref{eq:entropy} with $p = \sigma(\Delta)$. Then
\begin{equation}
  \frac{\partial \mathcal{H}}{\partial \Delta} \;=\; -\,\Delta\, p\,(1-p),
  \label{eq:dHdD}
\end{equation}
so the contribution of the term $-\beta \bar{\mathcal{H}}$ in \eqref{eq:loss} to the gradient
with respect to $\Delta_t$ is $+\beta\,\Delta_t\, p_t(1-p_t)$, and one gradient-descent step
moves
\begin{equation}
  \Delta_t \;\leftarrow\; \Delta_t \;-\; \eta\,\beta\,\Delta_t\,p_t(1-p_t) .
  \label{eq:restoring}
\end{equation}
This is a restoring force toward $\Delta_t = 0$ whose magnitude
$\beta|\Delta|\,\sigma(\Delta)(1-\sigma(\Delta))$ attains its maximum $0.2239\,\beta$ at
$|\Delta| = 1.543$ and decays as $\beta|\Delta|e^{-|\Delta|}$ thereafter.
\end{proposition}

\begin{proof}
$\mathrm{d}\mathcal{H}/\mathrm{d}p = \log\!\bigl((1-p)/p\bigr) = -\Delta$ and
$\mathrm{d}p/\mathrm{d}\Delta = p(1-p)$; the chain rule gives \eqref{eq:dHdD}. The stationary
point of $|\Delta|\sigma(\Delta)(1-\sigma(\Delta))$ is found numerically.
\end{proof}

\begin{table}[h]\centering
\begin{tabular}{rrr}
\toprule
$|\Delta|$ & $|\Delta|\,p(1-p)$ & $\mathcal{H}$ \\
\midrule
0 & 0.000 & 0.693 \\
1 & 0.197 & 0.582 \\
2 & 0.210 & 0.365 \\
4 & 0.071 & 0.090 \\
6 & 0.015 & 0.017 \\
8 & 0.003 & 0.003 \\
\bottomrule
\end{tabular}
\caption{The entropy force against the logit gap. Beyond $|\Delta|\approx 4$ it is
negligible.}
\end{table}

\noindent Three consequences follow, and they answer the question of what $\beta$ controls
\emph{inside} the network.

\begin{enumerate}
  \item \textbf{It is not a regulariser on the dynamics.} By \eqref{eq:delta},
        $\partial \mathcal{H}_t/\partial h_t = (\partial\mathcal{H}/\partial\Delta)\, w_\Delta$.
        The entropy gradient is parallel to $w_\Delta$ for every $t$. Any component of $h_t$
        orthogonal to $w_\Delta$ --- which is $N-1 = 127$ of the $128$ dimensions --- is
        invisible to it. A network can carry rich, structured dynamics at
        $\mathcal{H}\approx 0$, or none at all at $\mathcal{H} = \log 2$.
  \item \textbf{What it controls is exploration.} $\mathcal{H}$ high means $p_t$ stays away
        from $\{0,1\}$, so the agent still emits licks at times its current policy disfavours.
        Those are the only samples from which $\mathcal{L}_\pi$ can learn about actions it
        currently avoids. In \eqref{eq:lpi}, an action never taken contributes nothing.
  \item \textbf{$\beta_{\!f}$ is not a floor on $\mathcal{H}$.} $\beta$ is a coefficient on a
        force that, by Proposition~\ref{prop:ent}, has already vanished by the time the policy
        is deterministic. Once $|\Delta_t| \gtrsim 6$ the network is free to keep increasing
        $|\Delta_t|$ at essentially no entropy cost. This is the mechanism behind the observed
        collapse of variant H to $\bar{\mathcal{H}} = 0.005$ and $p_t \approx 0$
        (\S\ref{sec:consts}), despite $\beta_{\!f} = 0.03$: the run passes through good
        behaviour at update 50 and is not pulled back.
\end{enumerate}

\begin{remark}[Enforcing an actual floor]
A floor on $\mathcal{H}$ requires a term whose gradient does \emph{not} vanish with
$p(1-p)$ --- for instance a hinge $\beta\,\max(0,\, \mathcal{H}_{\min} - \bar{\mathcal{H}})$
with a Lagrange multiplier adapted online, or a hard clamp $p_t \leftarrow
\mathrm{clip}(p_t, \epsilon, 1-\epsilon)$ applied before sampling. Neither is implemented.
\end{remark}

\section{Constants}\label{sec:consts}

\begin{table}[h]\centering
\begin{tabular}{llrl}
\toprule
Symbol & Code name & Value & Meaning \\
\midrule
$\mathrm{d}t$ & \texttt{dt} & $0.02$\,s & integration and bin width \\
$\alpha$ & \texttt{alpha} & $0.2$ & leak, $\mathrm{d}t/\tau$ with $\tau = 100$\,ms \\
$N$ & \texttt{hidden} & $128$ & recurrent units \\
$\sigma_{\!\xi}$ & \texttt{noise} & $0.05$ & private noise SD \\
$B$ & \texttt{n\_envs} & $16$ & parallel environment copies \\
$\eta$ & \texttt{lr} & $4\times10^{-3}$ & Adam learning rate \\
$\gamma$ & \texttt{gamma} & $1.0$ & return discount \\
$c_V$ & \texttt{value\_coef} & $0.5$ & weight on $\mathcal{L}_V$ \\
$\epsilon$ & \texttt{min\_action\_prob} & $0.02$ & floor under every action probability \\
$G$ & \texttt{grad\_clip} & $1.0$ & gradient-norm clip \\
$c_{\mathrm{scale}}$ & \texttt{reward\_scale} & $0.1$ & global reward rescaling \\
$\lambda$ & \texttt{iti\_readout\_penalty} & $0$ or $1$ & per-step ITI penalty on $p_t$ \\
\midrule
$R_{\max}$ & \texttt{reward} & $10.0$ & water, before discounting \\
$\kappa$ & \texttt{discount\_rate} & $0.5$\,s$^{-1}$ & $e^{-\kappa\varepsilon}$ on the water \\
$k_{\mathrm{time}}$ & \texttt{time\_penalty} & $-0.05$ & per step, all phases \\
$k_{\mathrm{early}}$ & \texttt{early\_penalty} & $-1.0$ & lick before the delay elapsed \\
$k_{\mathrm{miss}}$ & \texttt{miss\_penalty} & $-2.0$ & answer window expired \\
$k_{\mathrm{iti}}$ & \texttt{iti\_lick\_penalty} & $-0.05$ or $-2.0$ & per accepted ITI lick \\
$k_{\mathrm{nc}}$ & \texttt{no\_cue\_lick\_penalty} & $-0.1$ & lick on a catch trial \\
$k_{\mathrm{to}}$ & \texttt{iti\_timeout\_penalty} & $-10.0$ & trial abandoned in the ITI \\
$n_{\mathrm{ref}}$ & \texttt{lick\_refractory} & $0.05$\,s & minimum inter-lick interval \\
$n_{\mathrm{cue}}$ & \texttt{cue\_duration} & $0.6$\,s & cue audible duration \\
$g$ & \texttt{reward\_rate\_gain} & $1.0$ & across-trial value gain \\
\bottomrule
\end{tabular}
\caption{Values as set in \texttt{variants.BASE}. Where two values are listed, the variants of
\S\ref{sec:variants} differ.}
\end{table}

\section{The variants}\label{sec:variants}

\begin{table}[h]\centering
\begin{tabular}{clccl}
\toprule
Run & Change from \texttt{BASE} & $k_{\mathrm{iti}}$ & $\lambda$ & restart \eqref{eq:restart} \\
\midrule
F & control & $-0.05$ & $0$ & on \\
G & $40\times$ the per-lick price & $-2.0$ & $0$ & on \\
H & no restart, high price & $-2.0$ & $0$ & \textbf{off} \\
I & per-step readout penalty & $-0.05$ & $1$ & on \\
J & both & $-0.05$ & $1$ & \textbf{off} \\
\bottomrule
\end{tabular}
\end{table}

\noindent Result of $150$ updates at seed $0$ (means of the last three log points): F, G and I
never reach a cue at all; only H and J, which drop \eqref{eq:restart}, learn anything. H attains
cue-success $0.64$ at update $50$ and then collapses to never-licking by update $125$, for the
reason given in Proposition~\ref{prop:ent}.

\end{document}
